高被引用論文のオブソレッセンス：
高被引用論文はその被引用の累積的な効果により長期にわたり上位の順位を維持するものの、最終的には減衰が確認され、新しい累積世代の論文にその座を明け渡す現象を確認できた。この交代までの期間は非常に長く最長で50年以上であった。うち、オブソレッセンスが確認された論文集合では、既存の報告と比べ、ピーク年（出版後に最大被引用数に達する年）が長くなる傾向があり、これは高被引用論文でのみ見られた。

高被引用論文掲載誌の特徴：
高被引用論文が必ずしもトップクラスの高インパクトファクター誌に掲載されるわけではないことが明らかになった。つまり、実際にはよりインパクトファクターの高い雑誌が収録されているにもかかわらず、高被引用論文掲載誌においてそのような雑誌の出現を確認できなかった。
とはいえ、高被引用論文掲載誌には、他のグループに比べインパクトファクターが高い雑誌も見られた。
高被引用論文の雑誌への嗜好については、特定の雑誌の確認はできなかったが、KLDによる分析から他のグループに比べより集中的であることが明らかになった。

高被引用テーマと高被引用論文引用テーマ：
高被引用論文は、研究テーマは集中的であるが、異なるテーマの論文から引用されやすいことが明らかになった。高被引用論文は、研究テーマおよび引用テーマそれぞれを個別に見た場合の多様性は低いが、研究テーマ対引用テーマの転換という観点では転換の程度が高くなる。もしこれがオブソレッセンスに影響しているのだとすれば、様々な研究領域から興味を持たれるようなテーマの論文は長く引用され続けるのかもしれない。前述の雑誌の特徴に照らせば、高被引用論文掲載誌に、幅広い領域から興味を持たれるであろう基礎領域や学際的分野の雑誌が見られる一方で、医学や医工学といった特定の応用領域の雑誌が見られないことも、これを示唆している。

長期的な傾向：
本研究ではさまざまな指標を用いたが、経年的な傾向が見られたものは少なかった。
唯一、引用テーマ/被引用テーマ/引用テーマ-被引用テーマ対のエントロピーに上昇傾向が見られた。

累積的な被引用データの影響：
累積の影響期間を変動させた順位相関の分析の結果、累積的な被引用を用いることにより順位の変動が大きくなった。すなわち累積的な被引用データを用いることは、順位等の相対的な評価のダイナミクスを上昇させることにつながる。
